\documentclass[11pt]{article}
\usepackage[utf8]{inputenc}
\usepackage{amsmath,amssymb,amsthm}
\usepackage{geometry}
\geometry{margin=1in}

\newtheorem{theorem}{Theorem}[section]
\newtheorem{lemma}[theorem]{Lemma}
\newtheorem{corollary}[theorem]{Corollary}
\newtheorem{definition}[theorem]{Definition}
\newtheorem{remark}[theorem]{Remark}

\title{Zero Topological Entropy from Hypoellipticity of the Lie Derivative}
\author{}
\date{}

\begin{document}
\maketitle

\begin{abstract}
Let \(M\) be a compact connected smooth manifold and let \(X\) be a smooth vector field on \(M\). Assume that the Lie derivative operator \(L_X\) is hypoelliptic in the following strong global sense: whenever \(f\in C^\infty(M)\) and \(u\in \mathcal D'(M)\) satisfy \(L_X u=f\), then \(u\in C^\infty(M)\). We prove that the flow \(g_t\) generated by \(X\) has zero topological entropy. The proof rests on a theorem of Dolgopyat and Saghin which states that positive topological entropy forces the existence of a smooth right-hand side \(f\) for which the cohomological equation \(L_Xu=f\) has a non-smooth distributional solution.
\end{abstract}

\section{Introduction and statement of the result}

Let \(M\) be a compact connected smooth manifold without boundary. Let \(X\in \mathfrak X^\infty(M)\) be a smooth vector field and let \(g_t\colon M\to M\) be the flow generated by \(X\). We denote by \(\mathcal D'(M)\) the space of distributions on \(M\). The Lie derivative along \(X\) acts on \(\mathcal D'(M)\) by duality: for \(u\in\mathcal D'(M)\) and \(\varphi\in C^\infty(M)\),
\[
(L_X u)(\varphi) = -u(L_X\varphi).
\]
The topological entropy of the flow \(g_t\) is, by definition, the topological entropy of the time-one map \(g_1\). We write it as \(h_{\mathrm{top}}(g_1)\).

The purpose of this note is to prove the following result.

\begin{theorem}\label{thm:main}
Let \(M\) be a compact connected smooth manifold and let \(X\) be a smooth vector field on \(M\). Assume that \(X\) satisfies the following property:
\begin{quote}
If \(f\in C^\infty(M)\) and \(u\in\mathcal D'(M)\) satisfy \(L_X u=f\), then \(u\in C^\infty(M)\).
\end{quote}
Then the flow \(g_t\) generated by \(X\) has zero topological entropy.
\end{theorem}

The hypothesis in Theorem~\ref{thm:main} is a global form of hypoellipticity for the operator \(L_X\). It is strictly stronger than the usual local hypoellipticity: for a non-singular vector field, local smooth solutions always exist, but global smooth solutions may fail to exist because of periodic orbits or Diophantine conditions.

\section{The key theorem of Dolgopyat and Saghin}

The proof of Theorem~\ref{thm:main} is an immediate consequence of the following deep theorem.

\begin{theorem}[Dolgopyat--Saghin]\label{thm:DS}
Let \(M\) be a compact smooth manifold and let \(X\) be a \(C^\infty\) vector field on \(M\). If the flow of \(X\) has positive topological entropy, then there exist \(f\in C^\infty(M)\) and \(u\in\mathcal D'(M)\) such that
\[
L_X u=f
\]
and \(u\notin C^\infty(M)\).
\end{theorem}

\begin{proof}[Reference]
This is the main result of D.~Dolgopyat and R.~Saghin, \emph{Hypoellipticity of the cohomological equation and entropy}, J. Mod. Dyn. \textbf{10} (2016), 1--17 (see also arXiv:1506.04759). The theorem is stated there for flows on compact manifolds; the case of a vector field \(X\) is the flow generated by \(X\). The proof uses Pesin's entropy formula, the existence of unstable manifolds for ergodic invariant measures with positive entropy, and the failure of smoothness of the unstable foliation. It constructs a smooth function \(f\) whose distributional solution \(u\) loses regularity along stable directions.
\end{proof}

\section{Proof of Theorem~\ref{thm:main}}

Assume, for the sake of contradiction, that the flow \(g_t\) has positive topological entropy:
\[
h_{\mathrm{top}}(g_1)>0.
\]
By the variational principle for the time-one map \(g_1\), there exists an ergodic \(g_1\)-invariant Borel probability measure \(\mu\) on \(M\) such that
\[
h_\mu(g_1)>0.
\]
Since \(g_1\) is the time-one map of the flow, \(\mu\) is also invariant under the whole flow \(g_t\).

Recall that the Lie derivative of the measure \(\mu\) is defined by
\[
(L_X\mu)(\varphi)=-\int_M L_X\varphi\,d\mu
\]
for every \(\varphi\in C^\infty(M)\). Because \(\mu\) is flow-invariant,
\[
\int_M L_X\varphi\,d\mu
=\int_M \frac{d}{dt}\Big|_{t=0}\varphi(g_t x)\,d\mu(x)=0.
\]
Hence \(L_X\mu=0\) in \(\mathcal D'(M)\).

Now apply the hypothesis of Theorem~\ref{thm:main} with \(f=0\) (which is certainly \(C^\infty\)) and \(u=\mu\) (which is a distribution). We conclude that \(\mu\) is a \(C^\infty\) function. That is, there exists a smooth volume form \(d\mathrm{Vol}\) on \(M\) and a function \(\rho\in C^\infty(M)\) such that
\[
d\mu=\rho\,d\mathrm{Vol}.
\]
Thus \(\mu\) is a smooth ergodic invariant measure with positive entropy.

We now apply Theorem~\ref{thm:DS} to the flow \(g_t\). Since \(h_{\mathrm{top}}(g_1)>0\), there exist \(f\in C^\infty(M)\) and \(u\in\mathcal D'(M)\) such that
\[
L_X u=f
\]
and \(u\notin C^\infty(M)\). This directly contradicts the hypothesis of Theorem~\ref{thm:main}. Therefore the assumption \(h_{\mathrm{top}}(g_1)>0\) was false, and we must have
\[
h_{\mathrm{top}}(g_1)=0.
\]
This proves Theorem~\ref{thm:main}.

\begin{remark}
The argument above only uses the hypothesis for the particular right-hand side \(f=0\) to conclude that every invariant measure is smooth, and then invokes the Dolgopyat--Saghin theorem for an ergodic invariant measure with positive entropy. In fact, the full strength of the hypothesis is not needed; it is enough to know that every invariant distribution is smooth and that the flow has an ergodic invariant measure with positive entropy whose distributional solution to the cohomological equation is not smooth. The latter is precisely the content of the Dolgopyat--Saghin theorem.
\end{remark}

\begin{thebibliography}{9}
\bibitem{DolgopyatSaghin}
D.~Dolgopyat and R.~Saghin,
\emph{Hypoellipticity of the cohomological equation and entropy},
J. Mod. Dyn. \textbf{10} (2016), 1--17.
arXiv:1506.04759.

\bibitem{Katok}
A.~Katok,
\emph{Lyapunov exponents, entropy and periodic orbits for diffeomorphisms},
Inst. Hautes Études Sci. Publ. Math. \textbf{51} (1980), 137--173.

\bibitem{BP}
L.~Barreira and Y.~Pesin,
\emph{Lyapunov exponents and smooth ergodic theory},
Univ. Lecture Ser. \textbf{23}, Amer. Math. Soc., Providence, RI, 2002.
\end{thebibliography}

\end{document}